\documentclass[11pt,a4paper]{article}
\usepackage[utf8]{inputenc}
\usepackage[T1]{fontenc}
\usepackage{natbib}
\usepackage[left=1in,right=1in,top=1in,bottom=1in]{geometry}
\usepackage{hyperref}

\title{Research Report: Agentic Science in Observational Cosmology}
\author{Cail M. Daley}
\date{March 2026}

\begin{document}

\maketitle

\section*{CMB Lensing and SPT-3G (2019--2024)}

My doctoral research at the University of Illinois focused on measurements of gravitational lensing of the cosmic microwave background (CMB) using data from the South Pole Telescope's third-generation receiver (SPT-3G).

CMB lensing---the subtle bending of CMB photons by the large-scale matter distribution---is a precise probe of structure growth, with a broad lensing kernel peaking at $z \sim 2$ that complements galaxy lensing at lower redshifts. The scientific challenge is extracting the lensing signal from CMB temperature and polarization maps with $\sim$arcminute resolution, where the signal is diffuse and degeneracies with other cosmological parameters are significant.

I led one of three independent lensing reconstruction pipelines, each extracting the lensing signal through different estimator choices to enable cross-validation. The reconstruction required simulation-based calibration of the estimator response across angular scales, and the covariance estimation used data-driven approaches to capture non-Gaussian correlations that analytical models missed. Cosmological parameters were extracted via Bayesian inference, marginalizing over instrumental and astrophysical nuisance parameters.

The dissertation delivered the deepest CMB lensing maps made to date as of 2024, with signal-to-noise ratio $S/N > 50$ and percent-level systematic uncertainties. This work was published in collaboration papers (Dutcher et al. 2021, Balkenhol et al. 2021, Pan et al. 2023, and the SPT-3G D1 data release papers in 2025--2026).

During this period, computational methods---simulations, high-dimensional parameter fitting, statistical estimator comparison---were essential, but the implementation was entirely human-driven.

\section*{Kinematic Lensing and Master's Supervision (2024--2025)}

At CosmoStat, I began supervising master's students on a new approach to gravitational lensing: measuring the distortions of galaxy velocity fields (kinematic lensing) rather than galaxy shapes (traditional cosmic shear). Kinematic lensing offers a fundamentally independent probe of the matter distribution, immune to the shape measurement systematics that limit cosmic shear.

The challenge is methodological: extracting lensing signals from spectroscopic survey data where the velocity measurements are noisy, redshifts are uncertain, and the distortions are subtle. I supervised three students on feasibility studies using simulated data and pilot analyses on real spectroscopic surveys (GAMA, DESI early data).

The methodology is statistical---developing velocity-based lensing estimators, simulating realistic survey noise, cross-validating on data sub-samples---but the later projects transitioned to agentic workflows, where students learned to design analyses and specify verification checks rather than writing every line of code themselves.

\section*{The UNIONS Cosmic Shear Analysis (2025--2026): Agentic Science}

The turning point came in mid-2025 when frontier language models became capable enough to autonomously implement substantial scientific analyses. I began the UNIONS cosmic shear analysis---a comprehensive measurement of weak lensing in the ultraviolet-near-infrared optical northern survey---intending to conduct it in the traditional human-implementation mode.

By the second month of work, it became clear that the cost-benefit had shifted. The agentic approach was faster, more reliable for routine implementation tasks, and---crucially---enabled a depth of verification and systematic exploration that would have been infeasible if I had implemented the code myself.

The UNIONS analysis involved:
\begin{itemize}
\item \textbf{Shape measurement calibration}: Developing photometric measurements of galaxy morphology, correcting for the point-spread function, and characterizing systematic biases in shape estimation. (\textit{~8,000 lines of code})
\item \textbf{Photometric redshift estimation}: Training and evaluating photometric redshift models, cross-validating against spectroscopic samples, and quantifying photo-z biases as a function of magnitude and color. (\textit{~4,000 lines})
\item \textbf{Two-point function measurement}: Computing correlation functions at multiple scales, validating against multiple independent codes (TreeCorr, NaMaster, Corrfuncs), and testing for systematic artifacts. (\textit{~6,000 lines})
\item \textbf{Covariance estimation}: Measuring sample covariance from simulations, computing model-based covariance predictions, and testing robustness to survey geometry and clustering assumptions. (\textit{~3,000 lines})
\item \textbf{Systematic modeling}: Implementing models for intrinsic alignments, baryonic feedback, and photometric errors; running multiple variants in parallel to assess sensitivity. (\textit{~15,000 lines})
\item \textbf{Bayesian inference}: Constructing likelihood functions, setting priors, running MCMC chains, and extracting cosmological constraints ($\Omega_m$, $\sigma_8$, intrinsic alignment amplitude, etc.). (\textit{~8,000 lines})
\item \textbf{Validation and visualization}: Generating hundreds of diagnostic plots, null tests on subsets of the data, consistency checks across magnitude and redshift bins. (\textit{~14,000 lines})
\end{itemize}

Total: \textasciitilde58,000 lines of Python, all verification plots, and the complete manuscript. All produced by AI agents under my supervision.

The shift in my role was qualitative. Instead of implementing code and debugging failures, I designed analyses and designed the tests that would tell me if something went wrong. I specified:
\begin{itemize}
\item Which estimators to use and why (TreeCorr for 2PCF, NaMaster for map-based correlations)
\item What null tests constituted evidence of systematic error (star-galaxy cross-correlations, consistency across redshift slices)
\item How to validate each component (comparison against prior analyses, scaling with survey area, robustness to parameter choices)
\item How much verification was enough (diminishing returns after \textasciitilde100 distinct tests)
\end{itemize}

The result is a cosmic shear analysis with verification and systematic exploration deeper than would have been feasible in a traditional 3-year postdoctoral timeline. The analysis measures weak lensing constraints on matter clustering and tests the $\Lambda$CDM model at percent-level precision.

This is the concrete evidence for the proposal's central claim: agentic science does not weaken rigor. It enables rigor at scales previously infeasible.

\section*{Diffusion: Adopting Agentic Workflows in Collaborations (2025--2026)}

In parallel with the UNIONS work, I have been helping other researchers adopt agentic workflows. Within the SPT collaboration, I have worked with several colleagues on adopting agentic methods for CMB power spectrum analyses and CMB lensing cross-correlations with optical surveys. The workflow is straightforward: design the analysis, specify the verification checks, and let agents handle implementation and iteration.

The practical experience has been consistent: end-to-end analyses that previously took 2--3 months of implementation now take 2--4 weeks, with more thorough testing because the cost of generating additional verification checks is low. The transition requires a genuine shift in how researchers approach their work---from implementing analyses to designing and specifying them---but colleagues who have made this shift report that it sharpens their scientific judgment, because the bottleneck moves from ``can I code this?'' to ``what should I check?''

\section*{Looking Forward}

The trajectory is clear: from human-implemented analysis with computational tools (SPT-3G), through a transitional period of supervised student work (kinematic lensing), to fully agentic scientific workflows producing publication-ready results (UNIONS). The next challenge is making this practice reproducible, verifiable, and teachable---not as a fixed toolkit but as an evolving methodology that adapts as models improve. The proposed research program addresses this through continued cosmological analyses (CMB lensing $\times$ Euclid, multi-probe weak lensing), formalization of verification practices, and training of early-career researchers who will carry these methods forward.

\end{document}
